Is the refinement loop needed, or would an exact solver on the atomic model do? We posed the 20 decision queries that produce the upper ends of the two-district brackets (for each state, party and rank the tightest margin at which the loop proved infeasibility; Table~\ref{tab:baseline}) to three exact methods under the same regime with a 300-second limit each: the CEGAR loop; CP-SAT on the atomic model of Section~\ref{sec:relax} (every precinct its own cell, exact parent-pointer connectivity and split row); and HiGHS on the independently written single-commodity-flow MILP of Section~\ref{sec:audit} on the atomic partition.

\begin{table}[!tbp]
\centering\footnotesize
\caption{Decision queries (``$q$ districts with share at least $m$?'', all answered ``no'' by the loop) posed to the CEGAR loop and to two monolithic exact solvers on the atomic model, 300 s limit, five CP-SAT workers (HiGHS: default threads). ``proof'' = proved infeasible; ``unknown'' = undecided at the limit. No solver returned a plan for a query that another proved infeasible.}\label{tab:baseline}
\adjustbox{max width=\linewidth}{\begin{tabular}{llcrrrr}
\toprule
State & Party & $q$ & margin $m$ & CEGAR & CP-SAT, atomic & HiGHS MILP, atomic \\
\midrule
Idaho & D & 1 & 50.0\% & proof ($<$1\,s) & proof (6\,s) & proof (2\,s) \\
Idaho & R & 1 & 76.0\% & proof (19\,s) & unknown ($>$300\,s) & unknown ($>$300\,s) \\
Idaho & R & 2 & 66.0\% & proof ($<$1\,s) & proof (1\,s) & proof (2\,s) \\
Maine & D & 1 & 63.0\% & proof ($<$1\,s) & proof (19\,s) & proof (41\,s) \\
Maine & D & 2 & 55.0\% & proof ($<$1\,s) & proof ($<$1\,s) & proof ($<$1\,s) \\
Maine & R & 1 & 55.0\% & proof ($<$1\,s) & proof (5\,s) & proof (2\,s) \\
Maine & R & 2 & 50.0\% & proof ($<$1\,s) & proof ($<$1\,s) & proof (1\,s) \\
Montana & D & 1 & 52.5\% & proof (91\,s) & unknown ($>$300\,s) & unknown ($>$300\,s) \\
Montana & D & 2 & 50.0\% & proof ($<$1\,s) & proof ($<$1\,s) & proof ($<$1\,s) \\
Montana & R & 1 & 69.5\% & proof (7\,s) & unknown ($>$300\,s) & unknown ($>$300\,s) \\
Montana & R & 2 & 58.5\% & proof ($<$1\,s) & proof (1\,s) & proof (1\,s) \\
New Hampshire & D & 1 & 57.5\% & proof ($<$1\,s) & unknown ($>$300\,s) & unknown ($>$300\,s) \\
New Hampshire & D & 2 & 54.0\% & proof ($<$1\,s) & proof ($<$1\,s) & proof ($<$1\,s) \\
New Hampshire & R & 1 & 50.0\% & proof ($<$1\,s) & unknown ($>$300\,s) & unknown ($>$300\,s) \\
Rhode Island & D & 1 & 70.5\% & proof (52\,s) & proof (19\,s) & unknown ($>$300\,s) \\
Rhode Island & D & 2 & 61.0\% & proof ($<$1\,s) & proof ($<$1\,s) & proof ($<$1\,s) \\
Rhode Island & R & 1 & 50.0\% & proof ($<$1\,s) & proof (1\,s) & proof ($<$1\,s) \\
West Virginia & D & 1 & 50.0\% & proof ($<$1\,s) & proof (11\,s) & proof (5\,s) \\
West Virginia & R & 1 & 77.0\% & proof (197\,s) & unknown ($>$300\,s) & unknown ($>$320\,s) \\
West Virginia & R & 2 & 70.0\% & proof ($<$1\,s) & proof (2\,s) & proof (7\,s) \\
\bottomrule
\end{tabular}
}
\end{table}

The loop decides all 20 (median 0.2 s; the slowest, West Virginia Republicans at 77\%, takes 197 s); atomic CP-SAT decides 14 and atomic HiGHS 13, and no two methods contradict each other, which is also a consistency check on the proofs (a plan returned where the loop proved infeasibility would refute the proof). The monolithic models are fast where the geography-free bound already nearly settles the query (every $q=2$ query and most zero-capacity queries: 0.3--41 s) and fail exactly on the seven best-district ($q=1$) queries near the transition margin, where the answer depends on how a county can be split. The largest contrast is New Hampshire: the county-level relaxation proves in about a tenth of a second (10 cells) that Democrats cannot reach 57.5\% and that Republicans cannot reach 50\% in any district, while neither monolithic solver decides either query in 300 s. The loop is not uniformly faster: atomic CP-SAT decides Rhode Island Democrats at 70.5\% in 19 s against 52 s for the loop. The comparison is between decision procedures under one time limit on one machine, not a statement about the best possible implementation of either.

\paragraph{Without the budget.}
The same test without a county-split budget ($s=\infty$) shows how much of the tractability comes from the budget. For three best-district queries---New Hampshire Democrats at 56\% and 58\%, Maine Democrats at 64\%---none of the three methods reaches a decision in 300 s: the loop has refined to 300--456 cells without a proof or a plan, and the atomic CP-SAT and HiGHS models return nothing, although a 40-second recombination run attains 57.0\% in New Hampshire, so the 56\% query is feasible. This agrees with the negative result: without the budget, certification of the interesting margins is out of reach for this method at precinct resolution, and it is the budget that turns the county level into an informative abstraction.
